\CapituloFinal{Conclusiones y recomendaciones}

\section*{Conclusiones}

El pensamiento computacional proporciona una metodología para transformar necesidades concretas en problemas que pueden ser abordados sistemáticamente. La descomposición, el reconocimiento de patrones, la abstracción y la generalización facilitan la construcción de modelos comprensibles y reutilizables.

El diseño algorítmico establece un vínculo entre la comprensión del problema y la escritura de código. El pseudocódigo, los diagramas de flujo y las pruebas de escritorio permiten evaluar la lógica antes de seleccionar detalles de implementación.

La programación convierte los modelos en soluciones ejecutables mediante datos, operaciones, estructuras de control y funciones. Sin embargo, la ejecución constituye solo una parte del proceso: las pruebas, la depuración y la mejora continua son esenciales para alcanzar resultados confiables y facilitar el mantenimiento.

El proyecto transversal de control de presupuesto muestra que una solución sencilla puede servir como entorno de práctica para todas estas competencias. La transferencia a problemas nuevos exige revisar sus condiciones particulares y no aplicar procedimientos de manera mecánica.

\section*{Recomendaciones}

\begin{enumerate}
\item Formular el problema y los criterios de aceptación antes de comenzar a programar.
\item Representar las soluciones con pseudocódigo y diagramas cuando ayuden a comunicar la lógica.
\item Practicar con problemas de dificultad creciente y explicar las decisiones de diseño.
\item Incluir datos normales, límites e inválidos en cada conjunto de pruebas.
\item Separar cálculos, validación y presentación mediante funciones con contratos claros.
\item Conservar registros de cambios y ejecutar pruebas de regresión.
\item Consultar documentación oficial y contrastar las referencias bibliográficas antes de publicar la obra.
\end{enumerate}
